% TOP_PROBLEM: {"id": 30005487, "title": "Campana–Peternell Conjecture for Fano Manifolds", "category_id": 5, "category": {"id": 5, "name": "algebraic_geometry", "display_name": "Algebraic Geometry", "description": "Geometric objects defined by polynomial equations.", "slug": "algebraic-geometry", "order_index": 5, "created_at": "2026-07-31T15:26:25.671Z"}, "status": "partially_solved"}
% FIRST_POSTED: 2026-09-27
% !TeX program = lualatex
% ATTEMPT_STATUS: unresolved
% Model: GPT-6 Astra Ultra; continuation dated 27 September 2026.
\documentclass[11pt]{article}
\usepackage[margin=25mm]{geometry}
\usepackage{fontspec}
\setmainfont{Latin Modern Roman}
\usepackage{amsmath,amssymb,mathtools}
\usepackage{unicode-math}
\setmathfont{Latin Modern Math}
\usepackage{xurl,hyperref}
\hypersetup{colorlinks=true,urlcolor=blue}
\setcounter{secnumdepth}{0}
\setlength{\parindent}{0pt}
\setlength{\parskip}{5pt}
\setlength{\emergencystretch}{3em}
\begin{document}
\section*{\texorpdfstring{Campana–Peternell Conjecture for Fano Manifolds}{Campana–Peternell Conjecture for Fano Manifolds}}\label{30005487}
\subsection{Definitions and mathematical statement}
A smooth complex projective variety $X$ is Fano when its anticanonical line bundle $K_X^{-1}$ is ample. A vector bundle $E$ is nef if the tautological quotient line bundle $\mathcal O_{{\mathbb{P}}(E)}(1)$ has nonnegative degree on every complete algebraic curve in ${\mathbb{P}}(E)$. The conjecture states
\[
 X\text{ Fano},\quad T_X\text{ nef}
 \quad\Longrightarrow\quad X\cong G/P,
\]
where $G$ is a connected complex semisimple algebraic group and $P$ a parabolic subgroup, so $G/P$ is projective and homogeneous. The isomorphism is algebraic. Nefness is a numerical condition weaker than ampleness; assuming an ample tangent bundle would substantially narrow the problem.
\par\smallskip\textit{Formulation and concept references:} \href{https://arxiv.org/abs/1407.6483}{[S1]}.

\subsection{Short English statement}
Must a smooth Fano variety with numerically nonnegative tangent bundle be a rational homogeneous space?
\subsection{Research attempt}

\paragraph{Scope and a source-status caution (27 September 2026).}
The hypothesis is nefness of the tautological quotient line bundle on
$\mathbb P(T_X)$, not ampleness of $T_X$ and not global generation.
Those stronger properties would remove major parts of the difficulty.
The supplied survey [S1] states the conjecture and its equivalent
global-generation formulation, and records established low-dimensional
and structural results. The elementary proofs below are reconstructions
of standard special cases and reductions, not proposed new classifications.

A more recent primary source requires particular care. Wang's
arXiv:2408.03799v2, dated 3 December 2025, retains an abstract claiming
that nefness of $T_X$ on a Fano manifold implies bigness and
semiampleness, but explicitly adds a notice that a gap in its proof
invalidates the main theorem's proof [E1]. The inspected version record
still lists v2. A promised corrected revision is not itself a proof.
This attempt does not use that claimed abundance result.

The approach is to understand restrictions to curves, prove the surface
case using the classical del Pezzo classification, and compare moving
rational curves with global vector fields. The final comparison produces
an exact missing implication rather than a false passage from numerical
positivity to symmetry.

\paragraph{Curve tests and elementary stability properties.}
We use the following form of the definition: a vector bundle $E$ on a
projective variety is nef if, for every morphism $f:C\to X$ from a
smooth projective curve and every quotient line bundle $f^*E\twoheadrightarrow Q$,
one has $\deg Q\geq0$. A quotient line defines a map to $\mathbb P(E)$
and pulls back its tautological bundle. Conversely, normalizing a curve
in $\mathbb P(E)$ gives such a quotient, which proves equivalence with
the definition in the statement.

Pullbacks and quotient bundles of nef bundles are nef, directly from
this criterion. A globally generated bundle is nef: every quotient line
on a curve is globally generated, and a nonzero section of a line bundle
has an effective zero divisor, so its degree cannot be negative.
These observations establish the easy direction for homogeneous spaces:
the infinitesimal action of a transitive algebraic group supplies global
vector fields spanning every tangent space.

The corresponding statement for subbundles is false. On $\mathbb P^1$
the exact sequence
\[
 0\longrightarrow\mathcal O(-1)\longrightarrow\mathcal O^{\oplus2}
 \longrightarrow\mathcal O(1)\longrightarrow0
\]
has a nef middle term and a non-nef subbundle. Thus in the relative
tangent sequence of a smooth morphism one may immediately use its
quotient, but may not deduce nefness of its kernel merely by inspection.
Any theorem asserting that fibres in the present setting again have
nef tangent bundle needs an additional argument.

\paragraph{Every rational curve has nonnegative splitting degrees.}
\textbf{Lemma 468.1.} If $T_X$ is nef and
$f:\mathbb P^1\to X$ is nonconstant, write
\[
 f^*T_X\cong\bigoplus_{i=1}^n\mathcal O(a_i)
\]
using the splitting theorem for bundles on $\mathbb P^1$. Then all
$a_i\geq0$, at least one $a_i\geq2$, and
$-K_X\cdot f_*[\mathbb P^1]\geq2$.

\emph{Proof.}
Each summand is a quotient line bundle of $f^*T_X$, so nefness gives
$a_i\geq0$. Since the ground field has characteristic zero and $f$
is nonconstant, its differential
$\mathcal O(2)=T_{\mathbb P^1}\to f^*T_X$ is not the zero map.
If every $a_i<2$, all its component homomorphisms would vanish because
$H^0(\mathbb P^1,\mathcal O(a_i-2))=0$. Thus one $a_i\geq2$.
Finally the anticanonical degree is the degree of $\det f^*T_X$,
namely $\sum_i a_i\geq2$.\hfill$\square$

This degree counts the parametrized curve with its mapping degree.
For a normalization map to an irreducible rational curve it gives the
ordinary anticanonical intersection. One cannot assume a splitting
of the special form $\mathcal O(2)\oplus\mathcal O(1)^a\oplus\mathcal O^b$
for every rational map: multiple covers already produce larger degrees.
For example a degree-$d$ map to $\mathbb P^1$ pulls its tangent bundle
back to $\mathcal O(2d)$.

The same splitting gives
\[
 H^1(\mathbb P^1,f^*T_X)=0,
 \qquad H^1(\mathbb P^1,f^*T_X(-p))=0.
\]
Under the standard deformation theory of morphisms, the first vanishing
makes the Hom-scheme smooth at $f$. The second says that evaluation
at $p$ is a submersion there: the differential is the evaluation map
from $H^0(f^*T_X)$, and its cokernel embeds into
$H^1(f^*T_X(-p))$ by the evaluation exact sequence.
Thus these curves move in families covering a neighbourhood of their
points. The deformation-theoretic interpretation is an established
input, also used in [S1, Lemma 2.8]; the displayed cohomology vanishing
is completely accounted for by the nonnegative splitting degrees.

\paragraph{Negative normal directions rule out many candidates.}
If $C\subset X$ is a smooth embedded curve, there is an exact sequence
\[
 0\longrightarrow T_C\longrightarrow T_X|_C
 \longrightarrow N_{C/X}\longrightarrow0.
\]
Nefness of $T_X$ forces $N_{C/X}$ to be nef as a quotient.
For a surface, $N_{C/X}$ is a line bundle of degree $C^2$, so
\[
 T_X\text{ nef}\quad\Longrightarrow\quad C^2\geq0
 \quad\text{for every smooth embedded curve }C.
 \tag{468.1}
\]
In particular a $(-1)$-curve on a surface is an immediate obstruction.
Its tangent bundle can have positive total degree and still fail
nefness: along an exceptional curve of a point blow-up the normal
quotient is $\mathcal O(-1)$, and in fact
\[
 T_X|_C\cong\mathcal O(2)\oplus\mathcal O(-1).
\]
The splitting follows because the extension class lies in
$H^1(\mathbb P^1,\mathcal O(3))=0$. Its determinant has degree one,
which makes clear that positivity of $-K_X$ on the curve does not
control every quotient of $T_X|_C$.

The same test rules out the tangent bundle of the blow-up of a smooth
$n$-fold at a point, for $n\geq2$. Its exceptional divisor is
$E\cong\mathbb P^{n-1}$ with normal line $\mathcal O_E(-1)$.
Restricting the quotient $T_X|_E\twoheadrightarrow N_{E/X}$ to a
line in $E$ gives a negative quotient of the restricted tangent bundle.
This statement concerns that specific blow-up geometry, not an assertion
that every birational morphism is a point blow-up.

\paragraph{The complete dimension-one and surface cases.}
In dimension one, a Fano curve has $\deg(-K)=2-2g>0$, so $g=0$
and it is $\mathbb P^1$. Its tangent line is $\mathcal O(2)$ and the
conclusion follows.

For dimension two we use the classical classification of smooth complex
del Pezzo surfaces as an explicit external input: each is $\mathbb P^2$,
$\mathbb P^1\times\mathbb P^1$, or the blow-up of $\mathbb P^2$ at
between one and eight points in the configurations giving an ample
anticanonical divisor. In the blow-up cases an exceptional $(-1)$-curve
exists. Equation (468.1) excludes every one of them. Hence a Fano surface
with nef tangent bundle is one of the first two surfaces.

Both remaining surfaces satisfy all hypotheses. The Euler sequence
\[
 0\longrightarrow\mathcal O_{\mathbb P^2}
 \longrightarrow\mathcal O_{\mathbb P^2}(1)^{\oplus3}
 \longrightarrow T_{\mathbb P^2}\longrightarrow0
\]
shows that $T_{\mathbb P^2}$ is globally generated and hence nef.
Its anticanonical line is $\mathcal O(3)$. On the product,
\[
 T_{\mathbb P^1\times\mathbb P^1}
 =\mathcal O(2,0)\oplus\mathcal O(0,2),\qquad
 -K=\mathcal O(2,2).
\]
Both summands are globally generated and the anticanonical line is
ample. These are homogeneous spaces for $\operatorname{PGL}_3$ and
$\operatorname{PGL}_2\times\operatorname{PGL}_2$, respectively.
This proof is a standard classification-based special case, not a
replacement for the much deeper results known in larger dimensions.

The Hirzebruch surfaces provide a useful independent numerical check.
For $\mathbb F_m=\mathbb P_{\mathbb P^1}(\mathcal O\oplus\mathcal O(m))$
with $m\geq0$, the minimal section $C_0$ has $C_0^2=-m$ and a
ruling fibre $F$ satisfies $F^2=0$, $C_0\cdot F=1$.
The canonical formula gives
\[
 -K_{\mathbb F_m}=2C_0+(m+2)F,\qquad
 -K\cdot C_0=2-m,\qquad -K\cdot F=2.
\]
Thus the Fano members are $m=0,1$, while nef tangent forces $m=0$
by the normal quotient along $C_0$. The Fano surface $\mathbb F_1$
is therefore a concrete example showing why the tangent hypothesis
adds real content beyond ampleness of the determinant.

\paragraph{Why enough global vector fields would finish the proof.}
\textbf{Proposition 468.2.} If $X$ is a smooth complex Fano variety
and $T_X$ is globally generated, then $X$ is a rational homogeneous
space $G/P$ with $G$ connected semisimple.

\emph{Proof.}
Let $G=\operatorname{Aut}^0(X)$. We use the standard identification
$\operatorname{Lie}(G)=H^0(X,T_X)$. The differential of the orbit map
$G\to X$, $g\mapsto g(x)$, at the identity is precisely evaluation
of global vector fields at $x$. Global generation therefore makes
every orbit open. Distinct orbits would partition the connected variety
into disjoint open subsets, so there is just one orbit.

For completeness, the Fano hypothesis removes any non-linear algebraic
group issue. Choose $m$ so that $-mK_X$ is very ample. Its canonical
functorial action embeds $\operatorname{Aut}(X)$ in the projective
linear group of its global sections: an automorphism acting trivially
on the embedded variety is the identity. Hence $G$ is a connected
linear algebraic group. Its connected solvable radical $R$ has a fixed
point on the complete variety $X$ by the Borel fixed-point theorem.
Because $R$ is normal and $G$ is transitive, it then fixes every point:
if $x$ is fixed, $r(gx)=g(g^{-1}rg)x=gx$.
The action is faithful, so $R$ is trivial. Thus $G$ is semisimple.
The stabilizer $P$ of a point has projective quotient $G/P\cong X$,
and is therefore parabolic.\hfill$\square$

The algebraic-group facts invoked here are standard: representability
of the automorphism group, the Lie-algebra identification, Borel's
fixed-point theorem and the projective-quotient characterization of
parabolic subgroups. This gives a complete reduction using those
facts; the unsolved part has not been hidden in a claim that all nef
bundles are generated by sections. The same reduction appears as the
global-generation version of the conjecture in [S1].

\paragraph{Local curve sections are not global vector fields.}
For a rational map $f:\mathbb P^1\to X$, nefness gives a surjective
evaluation map
\[
 H^0(\mathbb P^1,f^*T_X)\longrightarrow T_{X,f(p)}.
\]
To use Proposition 468.2 one instead needs
\[
 H^0(X,T_X)\longrightarrow T_{X,f(p)}
\]
to be onto. These spaces of sections are different. The restriction
map from global vector fields on $X$ to sections along $f$ has not
been shown surjective. Deforming one map $\mathbb P^1\to X$ changes
that curve; it does not produce an automorphism of the ambient variety.

There is no general numerical theorem repairing this distinction.
Even on a Fano surface a nef line bundle need not be globally generated.
Take a del Pezzo surface with $K_X^2=1$ and put $A=-K_X$.
It is ample and hence nef. Kodaira vanishing applied to
$A=K_X+2A$ gives $H^i(X,A)=0$ for $i>0$, while the same theorem
applied to $\mathcal O_X=K_X+A$ gives $\chi(\mathcal O_X)=1$.
Surface Riemann--Roch now gives
\[
 h^0(X,A)=\chi(A)=1+\tfrac12 A\cdot(A-K_X)=2.
\]
If this pencil were basepoint-free, it would define a morphism to
$\mathbb P^1$ with $A$ the pullback of $\mathcal O(1)$. Fibres over
two distinct points are disjoint, forcing $A^2=0$, contrary to
$A^2=1$. Thus $A$ is nef but not globally generated.
This is a line-bundle diagnostic on a Fano base, not a counterexample
involving that surface's tangent bundle; the latter is already ruled
out by its exceptional curves.

Likewise, semiampleness of $\mathcal O_{\mathbb P(T_X)}(1)$ would
give sections of a symmetric power after taking a multiple. It would
not, without an extra argument, yield the degree-one vector fields
needed to generate $T_X$. A proof of the weaker abundance assertion,
even if available, would still need a further homogeneity step.

\paragraph{Products and the limit of the reduction.}
Known examples are closed under products: if $X,Y$ have globally
generated tangent bundles, then
$T_{X\times Y}=p_X^*T_X\oplus p_Y^*T_Y$ is globally generated;
if both are Fano, the anticanonical line of the product is the exterior
tensor product of their anticanonical lines and is ample. If
$X=G/P$ and $Y=H/Q$, their product is
$(G\times H)/(P\times Q)$. This checks consistency of the predicted
class with the most elementary constructions.

It gives no decomposition theorem for an arbitrary Fano variety with
nef tangent bundle. Numerical splitting of intersection data is not
splitting of the variety, and an invariant family of rational curves
does not by itself construct the required algebraic group action.

\paragraph{Outcome and exact first gap.}
The attempt proves the dimension-one and surface cases using the stated
classical inputs, excludes negative normal quotients and point blow-ups,
derives the free-rational-curve vanishing, and proves that global
generation of $T_X$ gives the desired semisimple homogeneous description.
The examples test the differences between positive determinant, nefness,
generation along a curve and generation on the whole variety.

The first gap is the passage from all those curvewise nonnegative
quotients to enough global sections of $T_X$ at every point, or to
an independent construction of a transitive algebraic group action.
The supplied arguments establish neither passage in arbitrary dimension.
The main conjecture remains unresolved by this attempt.

\subsection{Sources}
\begin{itemize}
\item[S1]A survey on the Campana-Peternell Conjecture.
\url{https://arxiv.org/abs/1407.6483}.
\textit{Formulation source.}
\item[E1]Juanyong Wang, \emph{An abundance-type result for the tangent
bundles of smooth Fano varieties}, arXiv:2408.03799v2,
3 December 2025.
\url{https://arxiv.org/abs/2408.03799};
\url{https://arxiv.org/html/2408.03799v2}.
\textit{Version record and the author's opening gap notice inspected
27 September 2026. The affected main theorem is not used as an
established input.}
\item[E2]Roberto Mu\~noz, Gianluca Occhetta, Luis E. Sol\'a Conde,
Kiwamu Watanabe and Jaros\l aw A. Wi\'sniewski, full text of [S1],
revised 11 November 2015.
\url{https://arxiv.org/html/1407.6483}.
\textit{Conjecture 1.3, Proposition 2.1 and the rational-curve discussion
in Section 2.3 inspected 27 September 2026. Standard classification,
deformation and algebraic-group inputs are identified in the body.}
\end{itemize}
\end{document}
